\section{A Mellin transform with exact cancellation at every integer exponent}
\label{mpr:sec:main}

The next family keeps the polylogarithm order complex.  This is essential:
an identity proved only at positive integer orders would not justify the
order derivatives that produce Stieltjes constants below.  The starting
ingredients are the classical Fermi--Dirac integral for the polylogarithm,
the Lerch integral, and Euler's beta integral
\cite{DLMFPolylog,mpr:DLMFLerch,mpr:DLMFbeta}.
The contribution developed here is a finite factored operator for all
rational kernels in the family, together with joint continuation through
the Mellin and polylogarithm singularities.  The full-line family and its
operator were not located in the inspected manuscript or five newest
incoming reports.  This is a statement about the inspected corpus, not a
claim of worldwide priority.  Related work on integrals of products of
polylogarithms includes Freitas and Sofo
\cite{mpr:Freitas,mpr:Sofo}; their quoted finite-interval problems should
not be confused with the full-line transform proved here.

Let
\[
 \Omega=\mathbb C\setminus(-\infty,0],\qquad
 (\mathsf Ef)(s)=f(s-1),\qquad
 \mathscr Z(s)=s\zeta(s+1).
\]
The function $\mathscr Z$ is entire, with $\mathscr Z(0)=1$.
Products such as $(\mathsf E)_k$ are ordinary polynomials in the shift
operator, using the rising-factorial convention; in particular
$(\mathsf E)_2=\mathsf E(\mathsf E+1)$.
For $m\ge1$ define
\begin{equation}\label{mpr:eq:integral}
 \mathcal I_m(t,s;\lambda)
 =\int_0^\infty
       \frac{x^{t-1}\operatorname{Li}_s(-\lambda x)}{(1+x)^m}\,dx,
 \qquad -1<\Re t<m,\quad s\in\mathbb C,\quad\lambda\in\Omega.
\end{equation}
Here $x^{t-1}=\exp((t-1)\log x)$ uses the real logarithm, while the
polylogarithm uses its principal argument branch.  Thus its argument never
crosses $[1,\infty)$ for the displayed values of $\lambda$.
The stated $t$ strip is a common convergence domain for every order.
Special orders, notably negative integers, can have larger convergence
strips; no maximal-strip assertion for each individual order is intended.

\subsection{Analyticity of the ordinary integrals}

\begin{lemma}\label{mpr:lem:holomorphic}
The integral \eqref{mpr:eq:integral} is absolutely convergent and jointly
holomorphic in its three variables on the displayed domain.  Every fixed
number of derivatives with respect to $s$, $t$, and $\lambda$ may be taken
under the integral sign.  A $t$ derivative inserts a factor $\log x$, and
$\lambda\partial_\lambda$ replaces the order $s$ by $s-1$.
\end{lemma}

\begin{proof}
We record a compact-set bound that also covers negative and noninteger
orders.  If $s$ and $\lambda$ range over compact subsets of $\mathbb C$
and $\Omega$, respectively, then
\[
 \operatorname{Li}_s(-\lambda x)=O(x)\quad(x\downarrow0),\qquad
 \operatorname{Li}_s(-\lambda x)=O((1+\log x)^B)\quad(x\ge1)
\]
for a sufficiently large fixed $B$.  The first bound follows from the
normally convergent power series near zero.  For the second, choose an
integer $N$ with $\Re(s+N)>0$ throughout a slightly larger compact set and
use
\[
 \operatorname{Li}_{s+N}(-\lambda x)
 =-\frac1{\Gamma(s+N)}\int_0^\infty
        u^{s+N-1}\frac{\lambda x}{e^u+\lambda x}\,du.
\]
Apply $(x\partial_x)^N$.  For a complex variable on rays uniformly
separated from the negative real axis, every fixed logarithmic derivative
of $w/(1+w)$ is bounded by $C\min(|w|,1)$.  Splitting the $u$ integral
at $\log x+C_0$ therefore gives the asserted logarithmic bound.
The identity $(x\partial_x)^N\operatorname{Li}_{s+N}
=\operatorname{Li}_s$ holds first by the power series and then by analytic
continuation.  Cauchy's formula on the slightly larger compact parameter
sets supplies the same sort of bounds for all fixed parameter derivatives.

The resulting majorants at the two endpoints are
$x^{\Re t}$ and $x^{\Re t-m-1}(1+\log x)^B$.  They are locally uniformly
integrable throughout the common strip used in
\eqref{mpr:eq:integral}.  Inserting any fixed power of $\log x$ preserves
integrability.  The standard dominated holomorphy theorem completes the
proof.
\end{proof}

\subsection{The paired Lerch expression}

For $\Re a>0$, first on $|\lambda|<1$ within $\Omega$, put
\begin{equation}\label{mpr:eq:Hdefinition}
 \mathcal H_a(s,\lambda)
 =\operatorname{Li}_s(\lambda)
        -\lambda^a\Phi(\lambda,s,a),
 \qquad \lambda^a=\exp(a\operatorname{Log}\lambda).
\end{equation}
Its two terms must be continued together.  This convention has a concrete
ordinary-integral explanation.

\begin{lemma}\label{mpr:lem:H}
The function \eqref{mpr:eq:Hdefinition} extends jointly holomorphically
to $s\in\mathbb C$, $\Re a>0$, and $\lambda\in\Omega$.  For $\Re s>0$,
\begin{equation}\label{mpr:eq:Hkernel}
 \mathcal H_a(s,\lambda)
 =\frac1{\Gamma(s)}\int_0^\infty u^{s-1}
       \frac{w-w^a}{1-w}\,du,
 \qquad w=\lambda e^{-u},\quad
 w^a=\exp(a(\operatorname{Log}\lambda-u)).
\end{equation}
The apparent singularity at $w=1$ has the removable value $a-1$.
Consequently
\begin{align}
 \mathcal H_a(s,1)&=\zeta(s)-\zeta(s,a),
       \label{mpr:eq:Hone}\\
 \mathcal H_M(s,\lambda)&=\sum_{n=1}^{M-1}\frac{\lambda^n}{n^s}
       \quad(M=1,2,\ldots).
       \label{mpr:eq:Hinteger}
\end{align}
The pole of each zeta function in \eqref{mpr:eq:Hone} cancels.
\end{lemma}

\begin{proof}
For $|\lambda|<1$ and $\Re s>0$, the defining geometric series and the
gamma integral give \eqref{mpr:eq:Hkernel}.  The quotient is analytic at
$w=1$ because its numerator and denominator vanish to first order and
their ratio tends to $a-1$.  It and all fixed $u$ derivatives decay
exponentially as $u\to\infty$, locally uniformly in $a$ and $\lambda$.
They are bounded near $u=0$, including when $\lambda=1$.

To obtain the continuation in order without appealing to a divergent
Mellin integral, let $K_a(u,\lambda)=(w-w^a)/(1-w)$.  For any integer
$N\ge0$, integration by parts first on $\Re s>0$ gives
\[
 \mathcal H_a(s,\lambda)
 =\frac1{\Gamma(s+N)}\int_0^\infty
        u^{s+N-1}(-\partial_u)^N K_a(u,\lambda)\,du.
\]
The right side is holomorphic for $\Re s>-N$.  All boundary terms vanish
in the initial domain; the formulas agree on overlaps by continuation.
Since $N$ is arbitrary, they establish entire dependence on $s$ and the
claimed joint holomorphy.  Substitution of $\lambda=1$ proves
\eqref{mpr:eq:Hone} first where the separate zeta integrals converge and
then everywhere.  For a positive integer $M$, the kernel is the finite
geometric sum $w+\cdots+w^{M-1}$, proving
\eqref{mpr:eq:Hinteger}.
\end{proof}

\subsection{The full transform and the cancellation polynomial}

Define the explicit polynomial
\begin{equation}\label{mpr:eq:P}
 P_{m,t}(X)=(1-t-X)_{m-1}
           =\sum_{j=0}^{m-1}c_{m,j}(t)X^j.
\end{equation}

\begin{theorem}[Full Mellin transform]\label{mpr:thm:master}
For every parameter in \eqref{mpr:eq:integral},
\begin{equation}\label{mpr:eq:master}
 \boxed{\displaystyle
 \mathcal I_m(t,s;\lambda)
 =\frac{\pi}{(m-1)!\sin\pi t}
       \sum_{j=0}^{m-1}c_{m,j}(t)
                  \mathcal H_{m-t}(s-j,\lambda).}
\end{equation}
At $t=0,1,\ldots,m-1$ the numerator vanishes identically as a function
of $s$ and $\lambda$; the apparent poles are removable.  At $\lambda=1$
this becomes the entire-order Hurwitz formula
\begin{equation}\label{mpr:eq:hurwitz}
 \mathcal I_m(t,s;1)
 =\frac{\pi}{(m-1)!\sin\pi t}
       \sum_{j=0}^{m-1}c_{m,j}(t)
       \bigl(\zeta(s-j)-\zeta(s-j,m-t)\bigr).
\end{equation}
\end{theorem}

\begin{proof}
First take $m=1$, $0<\lambda<1$, $\Re s>0$, and $-1<\Re t<0$.
Insert the Fermi--Dirac integral in Lemma~\ref{mpr:lem:holomorphic}.
Absolute convergence permits interchanging the two integrations.  The
elementary partial fraction integral
\[
 \int_0^\infty\frac{x^t}{(1+x)(x+c)}\,dx
  =-\frac{\pi}{\sin\pi t}\frac{1-c^t}{c-1},\qquad c>1,
\]
follows from the beta integral.  With $c=e^u/\lambda$, it gives
\[
 \mathcal I_1(t,s;\lambda)
 =\frac\pi{\sin\pi t}
       \bigl(\operatorname{Li}_s(\lambda)
           -\lambda^{1-t}\Phi(\lambda,s,1-t)\bigr).
\]
The preceding two lemmas continue this identity in $s$ and $\lambda$,
and across the strip $-1<\Re t<1$ after removing $t=0$.

For $m\ge2$, integrate the derivative of
$x^t(1+x)^{1-m}\operatorname{Li}_s(-\lambda x)$ over the positive
axis.  On the common strip $-1<\Re t<m-1$ its boundary terms vanish.
Using $\mathcal I_m(t+1,s;\lambda)
=\mathcal I_{m-1}(t,s;\lambda)-\mathcal I_m(t,s;\lambda)$ yields
\begin{equation}\label{mpr:eq:recurrence}
 \mathcal I_m(t,s;\lambda)
 =\frac{m-1-t-\mathsf E}{m-1}\,
                     \mathcal I_{m-1}(t,s;\lambda).
\end{equation}
Iteration gives $P_{m,t}(\mathsf E)\mathcal I_1/(m-1)!$.

Initially on $|\lambda|<1$, the elementary Lerch shift relation gives
\[
 \mathcal H_{1-t}(s,\lambda)-\mathcal H_{m-t}(s,\lambda)
   =-\sum_{n=0}^{m-2}
       \frac{\lambda^{n+1-t}}{(n+1-t)^s}.
\]
For every summand, $\mathsf E$ acts as multiplication by $n+1-t$;
$P_{m,t}(n+1-t)=(-n)_{m-1}=0$.  Applying the polynomial therefore
annihilates the entire finite shift.  This proves
\eqref{mpr:eq:master} on a nonempty open strip, with the important
positive-real-part argument $m-t$.

At $t=k\in\{0,\ldots,m-1\}$, equation
\eqref{mpr:eq:Hinteger} says that its numerator equals
\[
 \sum_{n=1}^{m-k-1}\frac{\lambda^n}{n^s}
                    (1-k-n)_{m-1}=0.
\]
Every product in the last sum has a zero factor.  Thus all apparent
poles inside $-1<\Re t<m$ are removable.  The two sides are
holomorphic on that connected strip and agree on an open subset, proving
the full domain.  Equation \eqref{mpr:eq:Hone} gives the specialization
\eqref{mpr:eq:hurwitz}.
\end{proof}

For example, the base specialization is
\begin{equation}\label{mpr:eq:base}
 \int_0^\infty\frac{x^{t-1}\operatorname{Li}_s(-x)}{1+x}\,dx
  =\frac\pi{\sin\pi t}\bigl(\zeta(s)-\zeta(s,1-t)\bigr),
 \qquad -1<\Re t<1.
\end{equation}
The difference on the right is the generalized harmonic function
$H_{-t}^{(s)}$.  Its singularities at special orders have cancelled
before order differentiation is performed.

\section{The branch cancellation and a factored formula at resonance}
\label{mpr:sec:resonance}

For $|\lambda|<1$ in $\Omega$, define
\begin{equation}\label{mpr:eq:A}
 \mathcal A_s(\lambda)
    =\operatorname{Log}\lambda\,\operatorname{Li}_s(\lambda)
                         -s\operatorname{Li}_{s+1}(\lambda).
\end{equation}

\begin{theorem}[Continuation through the positive polylogarithm cut]
\label{mpr:thm:A}
The paired expression \eqref{mpr:eq:A} extends jointly holomorphically
to $s\in\mathbb C$, $\lambda\in\Omega$.  It satisfies
\begin{align}
 \mathcal A_s(\lambda)
    &=\int_0^\infty
        \frac{\operatorname{Li}_s(-\lambda x)}{x(1+x)}\,dx,
          \label{mpr:eq:Aintegral}\\
 \mathcal A_s(1)&=-\mathscr Z(s),
          \label{mpr:eq:Aone}\\
 \lambda\partial_\lambda\mathcal A_s(\lambda)
          &=\mathcal A_{s-1}(\lambda).
          \label{mpr:eq:Aderivative}
\end{align}
On the real interval $\lambda>1$, either consistent boundary value of
the two principal polylogarithms in \eqref{mpr:eq:A} gives the same
answer.  At $\lambda=1$ the full expression is removable even at
integer orders where its separate gamma and zeta expansions have poles.

If $u=\operatorname{Log}\lambda$ and $|u|<2\pi$, then
\begin{equation}\label{mpr:eq:localexpansion}
 \boxed{\displaystyle
 \mathcal A_s(e^u)
      =-\sum_{n=0}^\infty\mathscr Z(s-n)\frac{u^n}{n!}.}
\end{equation}
The series converges normally for $s$ in compact subsets of $\mathbb C$
and $|u|<2\pi$.  On the full $u$ disk it describes the continuation of
the germ at zero.  For the original principal composition one retains
$|\Im u|<\pi$, as is automatic when $u=\operatorname{Log}\lambda$.
\end{theorem}

\begin{proof}
Take the removable limit $t\to0$ in the $m=1$ transform, initially for
$|\lambda|<1$.  The Lerch derivative
$\partial_a\Phi(\lambda,s,a)=-s\Phi(\lambda,s+1,a)$ gives
\[
 \left.\partial_t
       \bigl(\operatorname{Li}_s(\lambda)
             -\lambda^{1-t}\Phi(\lambda,s,1-t)\bigr)\right|_{t=0}
 =\operatorname{Log}\lambda\,\operatorname{Li}_s(\lambda)
                         -s\operatorname{Li}_{s+1}(\lambda).
\]
This proves \eqref{mpr:eq:Aintegral}; the integral supplies the claimed
joint holomorphic continuation.  In particular, it proves that the two
boundary values across the positive cut agree, without assigning separate
divergent values.  For $\Re s>0$, taking $\lambda\to1$ from below in
\eqref{mpr:eq:A} gives $-s\zeta(s+1)$.  Both sides are entire in $s$,
proving \eqref{mpr:eq:Aone} at every order.  Differentiation of
\eqref{mpr:eq:Aintegral} proves \eqref{mpr:eq:Aderivative}.

There is also an explicit check of the branch cancellation.  For
noninteger $s$, the classical local polylogarithm expansion
\cite[Eq.~25.12.12]{DLMFPolylog} is
\[
 \operatorname{Li}_s(e^u)
  =\Gamma(1-s)(-u)^{s-1}
      +\sum_{n\ge0}\zeta(s-n)\frac{u^n}{n!},\qquad |u|<2\pi,
\]
on an initial compatible branch.  The two fractional-power terms in
$u\operatorname{Li}_s(e^u)-s\operatorname{Li}_{s+1}(e^u)$ cancel
because $\Gamma(1-s)=-s\Gamma(-s)$.  Combining the remaining series
gives precisely \eqref{mpr:eq:localexpansion}.

Every coefficient $\mathscr Z(s-n)$ is entire in $s$.  The zeta
functional equation and the gamma-ratio estimate imply, uniformly on a
compact set of orders,
\[
 \frac{|\mathscr Z(s-n)|}{n!}
        \le C_K n^{B_K}(2\pi)^{-n}
\]
for all sufficiently large $n$.  Consequently the series is normally
convergent on the claimed product domain.  Continuation in $s$ now
includes all integer orders and proves the claimed cancellation there.
The restriction on the original principal logarithm is stated explicitly
to avoid confusing a germ in $u$ with a globally single-valued expression
under the exponential map.
\end{proof}

\begin{theorem}[Every integer Mellin resonance]\label{mpr:thm:integer}
For $m\ge1$, $0\le k\le m-1$, $s\in\mathbb C$, and $\lambda\in\Omega$,
\begin{equation}\label{mpr:eq:integer}
 \boxed{\displaystyle
 \int_0^\infty\frac{x^{k-1}\operatorname{Li}_s(-\lambda x)}{(1+x)^m}\,dx
 =\frac{(\mathsf E)_k(1-\mathsf E)_{m-k-1}}{(m-1)!}
                                  \mathcal A_s(\lambda).}
\end{equation}
In particular,
\begin{equation}\label{mpr:eq:integerzeta}
 \int_0^\infty\frac{x^{k-1}\operatorname{Li}_s(-x)}{(1+x)^m}\,dx
 =-\frac{(\mathsf E)_k(1-\mathsf E)_{m-k-1}}{(m-1)!}
                                  \mathscr Z(s).
\end{equation}
The right side is a finite expression in entire functions.  Thus every
fixed number of order derivatives is valid, including at the values
where a displayed individual zeta factor has a pole.
\end{theorem}

\begin{proof}
At $t=0$, recurrence \eqref{mpr:eq:recurrence} and
\eqref{mpr:eq:Aintegral} give
$\mathcal I_m(0,s;\lambda)
=(1-\mathsf E)_{m-1}\mathcal A_s(\lambda)/(m-1)!$.
For $k\ge1$, multiplication of the kernel by
$x=(1+x)-1$ gives
\[
 \mathcal I_m(k,s;\lambda)
  =\mathcal I_{m-1}(k-1,s;\lambda)
                -\mathcal I_m(k-1,s;\lambda).
\]
Induct on $k$.  The claimed polynomials satisfy exactly this recurrence:
\begin{align*}
 &\frac{(X)_{k-1}(1-X)_{m-k-1}}{(m-2)!}
  -\frac{(X)_{k-1}(1-X)_{m-k}}{(m-1)!}\\
 &\hspace{15mm}
 =\frac{(X)_{k-1}(1-X)_{m-k-1}}{(m-1)!}
                   \bigl((m-1)-(m-k-X)\bigr)\\
 &\hspace{15mm}
 =\frac{(X)_k(1-X)_{m-k-1}}{(m-1)!}.
\end{align*}
This proves \eqref{mpr:eq:integer}.  Equation
\eqref{mpr:eq:Aone} proves \eqref{mpr:eq:integerzeta}.
\end{proof}

The first few kernels illustrate the factored operator:
\[
\begin{array}{c|c}
 (m,k)&(m-1)!^{-1}(X)_k(1-X)_{m-k-1}\\ \hline
 (1,0)&1\\
 (2,0)&1-X\\
 (2,1)&X\\
 (3,1)&(X-X^2)/2\\
 (3,2)&(X+X^2)/2\\
 (4,2)&(X-X^3)/6
\end{array}
\]
For instance, uniformly as an entire identity in $s$,
\begin{align}
 \int_0^\infty\frac{\operatorname{Li}_s(-x)}{(1+x)^2}\,dx
       &=-(s-1)\zeta(s),\label{mpr:eq:squared}\\
 \int_0^\infty\frac{x\operatorname{Li}_s(-x)}{(1+x)^4}\,dx
       &=\frac{(s-3)\zeta(s-2)-(s-1)\zeta(s)}6.
       \label{mpr:eq:quartic}
\end{align}
In \eqref{mpr:eq:squared}, the value at $s=1$ is $-1$, not zero;
in \eqref{mpr:eq:quartic}, the value of $(s-3)\zeta(s-2)$ at $s=3$
is $1$.  These are ordinary removable values of entire products.

\subsection{A bilinear identity isolating a single Stieltjes constant}

The same transform has a particularly short bilinear consequence.
One polylogarithm has a fixed negative integer order and is therefore
a rational function; the other order remains an arbitrary complex
variable.  No relation is differentiated outside its proven domain.

\begin{corollary}\label{mpr:cor:bilinear}
For every integer $d\ge1$, every $s\in\mathbb C$, and every
$\lambda\in\Omega$,
\begin{equation}\label{mpr:eq:bilinear}
 \boxed{\displaystyle
 \int_0^\infty\operatorname{Li}_s(-\lambda x)
                    \operatorname{Li}_{-d}(-x)\,\frac{dx}{x}
   =(-1)^d\mathcal A_{s-d}(\lambda).}
\end{equation}
The integral is ordinary and absolutely convergent, and is entire in
$s$.  At $\lambda=1$ it equals
$(-1)^{d+1}(s-d)\zeta(s-d+1)$, with the removable value understood.
Consequently, for every $r\ge1$,
\begin{equation}\label{mpr:eq:single-stieltjes}
 \boxed{\displaystyle
 \int_0^\infty\operatorname{Li}_{-d}(-x)
       \left.\partial_s^r\operatorname{Li}_s(-x)\right|_{s=d}
                 \frac{dx}{x}
        =(-1)^{d+r}r\gamma_{r-1}.}
\end{equation}
At derivative order zero the integral is $(-1)^{d+1}$.
\end{corollary}

\begin{proof}
Put $D=x\partial_x$ and $f(x)=1/(1+x)$.  Since
$f=1+\operatorname{Li}_0(-x)$, one has
$D^df=\operatorname{Li}_{-d}(-x)$ for $d\ge1$.
Every $D^jf$ is $O(x^{-1})$ at infinity.  At zero it is bounded
when $j=0$ and $O(x)$ when $j\ge1$.
Lemma~\ref{mpr:lem:holomorphic} therefore justifies $d$ successive
integrations by parts with respect to $dx/x$; every boundary term
vanishes.  They give
\[
 \int_0^\infty\operatorname{Li}_s(-\lambda x)D^df(x)\,\frac{dx}{x}
  =(-1)^d\int_0^\infty
       \frac{\operatorname{Li}_{s-d}(-\lambda x)}{x(1+x)}\,dx.
\]
Equation~\eqref{mpr:eq:Aintegral} proves \eqref{mpr:eq:bilinear}.
The same endpoint bounds give normal convergence in order, so the
order derivatives are valid.  At $\lambda=1$, use
$\mathcal A_v(1)=-\mathscr Z(v)$ and
$\mathscr Z^{(r)}(0)=r(-1)^{r-1}\gamma_{r-1}$ to obtain
\eqref{mpr:eq:single-stieltjes}.
\end{proof}

For example the choices $d=1$ and $d=2$ use only the kernels
\[
 \frac{\operatorname{Li}_{-1}(-x)}x=-\frac1{(1+x)^2},\qquad
 \frac{\operatorname{Li}_{-2}(-x)}x=\frac{x-1}{(1+x)^3}.
\]
Thus arbitrary Stieltjes indices occur in ordinary integrals with
fixed rational kernels.  The identity is a consequence of the
proved entire-order transform; no arithmetic independence assertion
is involved.

\section{All mixed logarithmic and Stieltjes identities in the family}
\label{mpr:sec:jets}

For $a\notin\{0,-1,-2,\ldots\}$ with the usual compatible branch,
write
\[
 \zeta(1+v,a)=\frac1v+
       \sum_{r=0}^\infty\frac{(-1)^r}{r!}\gamma_r(a)v^r,
 \qquad \gamma_r=\gamma_r(1).
\]
Only $\Re a>0$ is needed in the formulas below.

\begin{corollary}[Generalized Stieltjes differences]
\label{mpr:cor:stieltjes}
For $r\ge0$ and $0<\Re a<2$,
\begin{equation}\label{mpr:eq:generalizedstieltjes}
 \int_0^\infty\frac{x^{-a}}{1+x}
   \left.\partial_s^r\operatorname{Li}_s(-x)\right|_{s=1}\,dx
  =\frac{(-1)^r\pi}{\sin\pi a}
                     \bigl(\gamma_r-\gamma_r(a)\bigr).
\end{equation}
The value at $a=1$ is removable.  More explicitly it is
\begin{equation}\label{mpr:eq:aone}
 -\zeta^{(r)}(2)-r\zeta^{(r-1)}(2),
\end{equation}
where the second term is omitted for $r=0$.
\end{corollary}

\begin{proof}
Set $t=1-a$ in \eqref{mpr:eq:base} and differentiate the entire
difference $\zeta(s)-\zeta(s,a)$ $r$ times at $s=1$.
The common pole cancels before the coefficient is taken.  Lemma
\ref{mpr:lem:holomorphic} justifies the differentiated integral.
At $a=1$, use \eqref{mpr:eq:Aone} and differentiate $s\zeta(s+1)$ at
$s=1$, obtaining \eqref{mpr:eq:aone}.
\end{proof}

For a particularly explicit example, the half-argument distribution
formula $\zeta(s,1/2)=(2^s-1)\zeta(s)$ gives
\begin{equation}\label{mpr:eq:halfexample}
 \int_0^\infty\frac{x^{-1/2}}{1+x}
   \left.\partial_s\operatorname{Li}_s(-x)\right|_{s=1}\,dx
 =-\pi\bigl(2\gamma\log2+\log^2 2\bigr).
\end{equation}

\begin{corollary}[A single Stieltjes constant at every order]
\label{mpr:cor:quarticstieltjes}
For every $r\ge1$,
\begin{equation}\label{mpr:eq:quarticstieltjes}
 \boxed{\displaystyle
 \int_0^\infty\frac{x}{(1+x)^4}
    \left.\partial_s^r\operatorname{Li}_s(-x)\right|_{s=3}\,dx
 =\frac{r(-1)^{r-1}\gamma_{r-1}
       -2\zeta^{(r)}(3)-r\zeta^{(r-1)}(3)}6.}
\end{equation}
At order zero the same integral equals $(1-2\zeta(3))/6$.
\end{corollary}

\begin{proof}
The entire normalization satisfies
\[
 \mathscr Z(v)=1+\sum_{r\ge1}
                 \frac{(-1)^{r-1}\gamma_{r-1}}{(r-1)!}v^r,
 \qquad
 \mathscr Z^{(r)}(0)=r(-1)^{r-1}\gamma_{r-1}.
\]
Differentiate \eqref{mpr:eq:quartic} at $s=3$.  Its other entire factor
has derivative
$\mathscr Z^{(r)}(2)=2\zeta^{(r)}(3)+r\zeta^{(r-1)}(3)$.
This proves the formula, including the stated order-zero value.
\end{proof}

Thus the first two derivatives give the concrete ordinary integrals
\begin{align*}
 \int_0^\infty\frac{x}{(1+x)^4}
    \left.\partial_s\operatorname{Li}_s(-x)\right|_{s=3}\,dx
   &=\frac{\gamma-\zeta(3)-2\zeta'(3)}6,\\
 \int_0^\infty\frac{x}{(1+x)^4}
    \left.\partial_s^2\operatorname{Li}_s(-x)\right|_{s=3}\,dx
   &=-\frac{\gamma_1+\zeta'(3)+\zeta''(3)}3.
\end{align*}

\subsection{Every power of the logarithmic integration variable}

Write
\[
 \frac{\pi h}{\sin\pi h}=\sum_{a\ge0}b_{2a}h^{2a},\qquad
 b_0=1,\qquad b_{2a}=2(1-2^{1-2a})\zeta(2a)\quad(a\ge1).
\]
For example $b_2=\zeta(2)$ and $b_4=7\zeta(4)/4$.

\begin{theorem}[All logarithmic moments]\label{mpr:thm:logmoments}
For every integer $\ell\ge0$ and every complex $s$,
\begin{equation}\label{mpr:eq:logmoments}
 \boxed{\displaystyle
 \int_0^\infty
       \frac{\log^\ell x\operatorname{Li}_s(-x)}{x(1+x)}\,dx
 =-\ell!\sum_{a=0}^{\lfloor\ell/2\rfloor}
        b_{2a}\frac{(s)_{\ell+1-2a}}{(\ell+1-2a)!}
                    \zeta(s+\ell+1-2a).}
\end{equation}
Every product $(s)_n\zeta(s+n)$ occurring here is entire, with its
removable value understood.  Every fixed number of order derivatives
can be taken on both sides.
\end{theorem}

\begin{proof}
The Hurwitz Taylor series at $a=1$, obtained by repeated use of
$\partial_a\zeta(s,a)=-s\zeta(s+1,a)$, gives
\[
 \zeta(s)-\zeta(s,1-h)
    =-\sum_{n\ge1}\frac{(s)_n}{n!}\zeta(s+n)h^n,
       \qquad |h|<1.
\]
The putative pole at $s=1-n$ cancels against the corresponding linear
factor of $(s)_n$.  The identity holds normally for $s$ in compact sets
by the holomorphy of the Hurwitz difference and Cauchy's formula.
Multiply this series by $\pi/\sin\pi h$, take the coefficient of
$h^\ell$, and apply Lemma~\ref{mpr:lem:holomorphic} to
\eqref{mpr:eq:base}.  This proves the formula and its order derivatives.
\end{proof}

The first two nonconstant logarithmic moments are
\begin{align*}
 \int_0^\infty\frac{\log x\operatorname{Li}_s(-x)}{x(1+x)}\,dx
   &=-\frac{s(s+1)}2\zeta(s+2),\\
 \int_0^\infty\frac{\log^2x\operatorname{Li}_s(-x)}{x(1+x)}\,dx
   &=-\frac{s(s+1)(s+2)}3\zeta(s+3)
                           -2\zeta(2)s\zeta(s+1).
\end{align*}
In particular,
\[
 \int_0^\infty\frac{\log x}{x(1+x)}
      \left.\partial_s\operatorname{Li}_s(-x)\right|_{s=-1}\,dx
       =\frac{\gamma-1}{2}.
\]
This follows by expanding
$-s(s+1)\zeta(s+2)/2$ at $s=-1$; its value before differentiating is
$1/2$.

For a general kernel and general scale, all Mellin derivatives at a
resonant exponent have an equally finite coefficient prescription.
Let
\[
 N_m(t,s,\lambda)
     =\sum_{j=0}^{m-1}c_{m,j}(t)\mathcal H_{m-t}(s-j,\lambda).
\]
Since $N_m(k,s,\lambda)=0$, Theorem~\ref{mpr:thm:master} gives
\begin{equation}\label{mpr:eq:mixedcompiler}
 \left.\partial_t^\ell\partial_s^r
                      \mathcal I_m(t,s;\lambda)\right|_{t=k}
 =\frac{(-1)^k\ell!}{(m-1)!}
   [h^{\ell+1}]\frac{\pi h}{\sin\pi h}\,
                       \partial_s^r N_m(k+h,s,\lambda).
\end{equation}
Only finitely many derivatives and coefficients are required for each
fixed $\ell,r,m$.  This formula also retains mixed $s$ and $\lambda$
derivatives, since $\lambda\partial_\lambda$ simply shifts the order.
It specifies the transverse coordinate $t-k$ and never differentiates a
formula that was known only on the resonant set.

\section{Exact scaled examples and further directions}
\label{mpr:sec:examples}

Put $L=\log2$.  Classical values at $1/2$ yield
\begin{align}
 \mathcal A_1(\lambda)
       &=\operatorname{Li}_2(1-\lambda)-\zeta(2),
            \label{mpr:eq:A1}\\
 \mathcal A_2(1/2)
       &=-\frac74\zeta(3)+\frac{\pi^2L}{12}+
                                               \frac{L^3}{6},
            \label{mpr:eq:Ahalf}\\
 \mathcal A_2(2)
       &=-\frac74\zeta(3)-\frac{\pi^2L}{4}.
            \label{mpr:eq:Atwo}
\end{align}
Equation \eqref{mpr:eq:A1} is Euler's dilogarithm reflection identity in
the present normalization.  Equation \eqref{mpr:eq:Ahalf} follows from
the standard formulas for $\operatorname{Li}_2(1/2)$ and
$\operatorname{Li}_3(1/2)$.
For \eqref{mpr:eq:Atwo}, the compatible boundary values are
\[
 \operatorname{Li}_2(2)=\frac{\pi^2}{4}-i\pi L,\qquad
 \operatorname{Li}_3(2)=\frac78\zeta(3)+\frac{\pi^2L}{4}
                                           -\frac{i\pi L^2}{2}.
\]
They follow from classical inversion and the half-argument evaluations;
their imaginary parts cancel in $L\operatorname{Li}_2(2)
-2\operatorname{Li}_3(2)$.  The opposite compatible boundary convention
gives their conjugates and the same real answer.

Theorem~\ref{mpr:thm:integer} consequently gives, without any integer
relation search,
\begin{align}
 \int_0^\infty\frac{\operatorname{Li}_2(-x/2)}{x(1+x)}\,dx
       &=-\frac74\zeta(3)+\frac{\pi^2L}{12}+\frac{L^3}{6},
             \label{mpr:eq:ex1}\\
 \int_0^\infty\frac{\operatorname{Li}_2(-2x)}{x(1+x)}\,dx
       &=-\frac74\zeta(3)-\frac{\pi^2L}{4},
             \label{mpr:eq:ex2}\\
 \int_0^\infty\frac{\operatorname{Li}_3(-2x)}{(1+x)^3}\,dx
       &=\frac{\pi^2(1-L)}8-\frac78\zeta(3).
             \label{mpr:eq:ex3}
\end{align}
All are ordinary absolutely convergent integrals.  The last example
uses the short operator $(\mathsf E-\mathsf E^2)/2$ and
$\mathcal A_1(2)=-\pi^2/4$.

The results suggest several precise continuations.
\begin{enumerate}
\item Determine which noninteger beta weights retain a finite expression
in shifts of $\mathscr Z(s)$.  The present factored operator uses the
integrality of the denominator exponent, and does not settle finite
closure for a general complex exponent.
\item Extend the factored operator to kernels with several independent
negative poles.  Partial fractions give a finite reduction to scaled
instances of $\mathcal A_s$; an explicit confluent formula should keep
coalescing poles and order derivatives regular simultaneously.
\item Determine which bilinear full-line integrals with two independent
polylogarithm orders admit finite reductions of comparable simplicity.
A product introduces an additional convolution; the present one-factor
theorem does not settle that question.
\item Study the analytic continuation of the local series
\eqref{mpr:eq:localexpansion} around $u=\pm2\pi i$ and its monodromy.
The theorem proves normal convergence within the stated disk, but does
not supply a global single-valued continuation in $u$.
\item Compare the precise all-order transform and resonance formulas with
the older Mellin-transform tables and subsequent Euler-sum literature
before asserting priority beyond the inspected corpus.
\end{enumerate}
